% =====================================================================
%  Seksjon: Kontroller-blokkering -- hvorfor ROSCO hindrer at en
%           nett-/elektrisk forstyrrelse eksiterer tarnmodene (side-to-side
%           OG fore-aft) i den hoy-fidelitets OpenFAST+ROSCO-turbinen.
%           Samler bevis-testene og forklarer blokkeringen mode-for-mode.
%  Kilder:  docs/openfast_excitation.tex        (detaljert SS-behandling)
%           docs/STATUS_OG_VEIVALG.md, kap. 2-3  (testtabell + rotarsak)
%           test1002/ControlData/ROSCO.IEA15MW.IN (VSContrl, OL_Mode, TRA_Mode,
%                                                   F_LPFCornerFreq)
%  Preamble: amsmath, amssymb, graphicx, booktabs, siunitx, hyperref
% =====================================================================

\section{Controller Blocking: Why ROSCO Prevents Electrical Excitation of the Tower Modes}
\label{sec:rosco-blocking}

This project studies electro-mechanical interactions in offshore wind turbines
coupled to a grid --- here an IEA \SI{15}{\mega\watt} turbine connected to the
LEOGO islanded platform network. One such interaction is whether an
electrical/grid disturbance can reach the turbine's structural tower modes. In
the high-fidelity OpenFAST turbine that particular path fails: neither the
side-to-side (SS) nor the fore-aft (FA) tower mode can be excited through any
electrical input, because the ROSCO controller owns the one channel that would
carry the disturbance into the structure. This section gathers
the tests that established this negative result --- the tests that ``proved'' the
blocking --- for \emph{both} tower modes, and explains mode-by-mode why the
electrical path is closed. It is a pivot point of the methodology: it is precisely
this blocking that motivates the ROSCO open-loop torque injection that bypasses it
(\S\ref{subsec:se-openfast}) and the reduced \texttt{WindTurbineTower} whose
electrical-to-mechanical path is open by construction
(\S\ref{sec:simplified-tower-model}). The detailed SS experiment log lives in
\S\ref{sec:openfast-excitation}; here the focus is the evidence and the causal
argument.

\subsection{The two candidate drive paths}
\label{subsec:rb-paths}

For a grid disturbance to reach a tower mode it must modulate one of two
mechanical loads: the \emph{generator torque}, whose nacelle reaction moment rolls
the nacelle laterally and bends the tower side to side; or the \emph{rotor
thrust}, which bends the tower fore and aft. The two first tower modes are nearly
degenerate in frequency
(\(f_{\mathrm{ss}}\approx f_{\mathrm{fa}}\approx\SI{0.234}{\hertz}\)), so what
decides whether an electrical disturbance can reach each of them is not the
frequency but whether the disturbance can act on torque (for SS) or on thrust
(for FA). Table~\ref{tab:rb-paths} states the mapping that the tests below probe.

\begin{table}[t]
  \centering
  \caption{The two electrical-to-structure drive paths and the controller that
  guards each in the OpenFAST-FMU (\texttt{VSContrl}=5).}
  \label{tab:rb-paths}
  \begin{tabular}{lll}
    \toprule
    Tower mode & Mechanical drive & Grid access in the FMU \\
    \midrule
    Side-to-side & generator torque \(T_e\) & only via ROSCO's torque command \\
    Fore-aft     & rotor thrust \(F_T\)      & no electrical thrust channel at all \\
    \bottomrule
  \end{tabular}
\end{table}

That the torque-to-tower coupling itself \emph{exists} in the model is not in
doubt: ROSCO even ships a Tower Resonance Avoidance mode (\texttt{TRA\_Mode}) that
deliberately modulates \emph{generator torque} to steer the rotor away from a
tower frequency. The coupling is real; the question the tests answer is whether an
\emph{external, electrically imposed} torque or thrust can use it.

\subsection{Tests that closed the torque path (side-to-side)}
\label{subsec:rb-torque}

Three tests, each a baseline-versus-modulated comparison, established that the
generator-torque channel is closed to external commands.

\begin{enumerate}
  \item \textbf{Generator-torque command} (\texttt{GenSpdOrTrq}). A commanded
        torque modulation at the SS frequency left the simulated generator torque
        \emph{bit-identical} to the unmodulated baseline --- the command has no
        effect whatsoever.
  \item \textbf{Electrical-power command} (\texttt{ElecPwrCom},
        \(\pm\SI{20}{\percent}\) at \SI{0.236}{\hertz}). Ignored identically in
        Region~2 (\SI{8}{\metre\per\second}) and Region~3: the \texttt{GenTq}
        ripple and the SS acceleration envelope were indistinguishable from
        baseline (\S\ref{subsec:of-electrical}, Table therein).
  \item \textbf{Opening the external-torque path} (\texttt{VSContrl}=4). The one
        configuration that would let an external torque become the actual
        generator torque \emph{crashes} the FMU at the first \texttt{doStep}
        (access-violation in the FMI call); the binary lacks the
        Simulink/Labview pointer hookup and it is not fixable from configuration
        (\S\ref{subsec:of-vscontrl4}).
\end{enumerate}

The cause is architectural. With \(\texttt{VSContrl}=5\) ROSCO computes the
generator torque \emph{internally} from the measured generator speed --- in
Region~2 the optimal law
\begin{equation}
  T_{\mathrm{gen}} = K_{\mathrm{opt}}\,\omega_g^2,
  \label{eq:rb-komega2}
\end{equation}
and does not read the external \texttt{GenSpdOrTrq}/\texttt{ElecPwrCom} ports. The
small torque ripple seen in the runs is ROSCO's own speed-reaction to the start-up
transient, not a response to the command. Table~\ref{tab:rb-evidence} collects the
evidence. The verdict for the torque path, and hence for the SS mode, is
unambiguous: \emph{it is closed}.

\begin{table}[t]
  \centering
  \caption{The controller-blocking evidence (OpenFAST-FMU, \texttt{VSContrl}=5).
  Every electrical route into the generator torque is either ignored or crashes.}
  \label{tab:rb-evidence}
  \begin{tabular}{lll}
    \toprule
    Test & Result & Root cause \\
    \midrule
    Torque cmd.\ \texttt{GenSpdOrTrq}      & bit-identical \texttt{GenTq}        & ROSCO owns torque (\ref{eq:rb-komega2}) \\
    Power cmd.\ \texttt{ElecPwrCom} \(\pm20\%\) & ignored, Region~2 \emph{and} 3 & same --- torque set internally \\
    \texttt{VSContrl}=4 (ext.\ torque)     & crash at first \texttt{doStep}      & missing Simulink pointer; needs rebuild \\
    ROSCO filters hiding SS?               & no (see \S\ref{subsec:rb-filters})  & filters shape response, not access \\
    \bottomrule
  \end{tabular}
\end{table}

\subsection{Ruling out the ROSCO filters as the cause}
\label{subsec:rb-filters}

A natural suspicion is that the blocking is really a \emph{filtering} effect ---
that ROSCO's low-pass and notch filters (\texttt{F\_LPFCornerFreq} and the tower
notch used by the SS-damping / resonance-avoidance logic) are actively cancelling
the SS excitation before it reaches the structure. This was tested and rejected.
Those filters act only on the \emph{speed signal that ROSCO measures}, shaping
\emph{how ROSCO responds}; they are not the reason the network fails to reach the
SS mode. The disturbance is stopped one step earlier, at the torque command
itself: because \(\texttt{VSContrl}=5\) makes the torque a pure function of
measured speed \eqref{eq:rb-komega2}, the external command never enters the torque
in the first place, so there is nothing for a notch filter to remove. The blocking
is architectural, not spectral --- an important distinction, because it means no
filter re-tuning could open the path; only a different converter interface could.

\subsection{The fore-aft mode is doubly unreachable}
\label{subsec:rb-foreaft}

The fore-aft mode is blocked even more firmly than the side-to-side mode, for a
different reason. FA is driven by rotor \emph{thrust}, and the FMU exposes no
electrical input that acts on thrust at all --- there is simply no
grid-to-thrust channel to command. The only conceivable route is the indirect one
already written in \eqref{eq:se-fa-path},
\[
  \text{grid}\to P_e\to T_e\to \omega_m\to \lambda\to C_T\to F_T
  \to q_{\mathrm{fa}},
\]
and every gate on that route is shut: its first link is the very torque channel
that ROSCO owns (\S\ref{subsec:rb-torque}), so \(T_e\) barely moves; and even if it
did, the huge rotor inertia would keep \(\omega_m\), and therefore \(\lambda\),
\(C_T\) and the thrust, nearly constant at fixed wind. The FA mode is thus blocked
twice over --- no direct electrical thrust channel, and the indirect
torque\(\to\)speed\(\to\)thrust channel throttled by both ROSCO and the rotor
inertia. Neither tower mode, then, can be electrically excited in the FMU.

This is also why the FA/SS distinction survives into the reduced model, where the
torque path \emph{is} open: there SS is driven and FA is not, by a factor of
\(\sim\!30\) at resonance (\S\ref{subsec:se-selectivity},
\S\ref{subsec:nt-resonance}). The controller blocking in the FMU and the drive-path
selectivity in the reduced model are two views of the same fact --- the grid can
only ever reach the tower through torque, and only the side-to-side mode listens to
torque.

\subsection{The core: full-converter decoupling}
\label{subsec:rb-core}

Putting the pieces together, the OpenFAST-FMU behaves as a \emph{full-converter
decoupling with no return path}. With \(\texttt{VSContrl}=5\) the OpenFAST
generator torque is entirely ROSCO's; the network disturbance reaches only the
co-simulation \emph{wrapper} generator mass --- which is why it \emph{does} excite
the \SI{3.49}{\hertz} drivetrain torsional mode of \S\ref{sec:forcedresponse} ---
but never OpenFAST's own structural solver, and hence never the tower modes. The
electrical and structural sides of the high-fidelity model are separated by the
controller, and there is no path back.

Two design choices follow directly, and both are pursued elsewhere in this report.
To prove that the OpenFAST \emph{tower} really would amplify a torque oscillation
if one could be imposed, the torque is injected through ROSCO's own open-loop
control table, bypassing the \(K\omega^2\) law entirely
(\S\ref{subsec:se-openfast}); that experiment recovers the sharp SS resonance and
so isolates the structural selectivity from the controller blocking. And to obtain
a genuinely \emph{closed} electrical drive of the tower --- a real grid
disturbance producing the torque --- the study moves to the reduced
\texttt{WindTurbineTower} (\S\ref{sec:simplified-tower-model}), where
\(T_e = P_e/(\omega_{e,\mathrm{filt}}\eta)\) enters the modal equation explicitly
and no controller intercepts it. The blocking result of this section is thus not a
dead end but the hinge of the whole methodology.

\subsection*{Sources}
\label{subsec:rb-sources}

\begin{itemize}
  \item ROSCO controller (\texttt{VSContrl} torque law, open-loop control,
        \texttt{TRA\_Mode} tower resonance avoidance, \texttt{F\_LPFCornerFreq}):
        Abbas et al.\ (2022), \emph{A reference open-source controller for fixed
        and floating offshore wind turbines}, Wind Energy Science 7, 53--73,
        \href{https://doi.org/10.5194/wes-7-53-2022}{doi.org/10.5194/wes-7-53-2022};
        documentation \href{https://rosco.readthedocs.io/}{rosco.readthedocs.io};
        source \href{https://github.com/NREL/ROSCO}{github.com/NREL/ROSCO}.
  \item OpenFAST / ServoDyn (\texttt{VSContrl} options, FMU interface):
        \href{https://openfast.readthedocs.io/}{openfast.readthedocs.io}.
\end{itemize}
